Many recent successful generative models have their origin in physical processes, including diffusion models~\cite{Ho2020DenoisingDP, Song2020ScoreBasedGM, Karras2022ElucidatingTD} and Poisson flow generative models~\cite{Xu2022PoissonFG, Xu2023PFGMUT}. These models involve a forward process that transforms the data distribution into a chosen smooth distribution, and a backward process that iteratively reverses the forward process. For instance, in diffusion models, the forward process is the diffusion process with no learned parameters:
\begin{equation*}
%\label{eq:sgm_forward}
    \mathrm{d}x = \sqrt{2\dot{\sigma}(t)\sigma(t)}\mathrm{d}W_t,
\end{equation*}
where $\sigma(t)$ is a predefined noise schedule increasing with $t$, and $W_t\in\mathbb{R}^d$ is the standard Wiener process. For simplicity, we omit an additional scaling function for other variants of diffusion models as in EDM~\cite{Karras2022ElucidatingTD}. Under this notation, the marginal distribution at time $t$ is the convolution of data distribution $p_0=p_{\textrm{data}}$ and a Gaussian kernel, \ie $p_t = p_0 * \gN(\bm{0}, \sigma^2(t)\bm{I}_{d\times d})$. The prior distribution is set to $\gN(\bm{0}, \sigma^2(T)\bm{I}_{d\times d}) $ since $p_T$ is approximately Gaussian with a sufficiently large $T$. Sampling of diffusion models is done via a reverse-time SDE~\citep{anderson1982reverse} or a marginally-equivalent ODE~\cite{Song2020ScoreBasedGM}:
\begin{align*}
    &\text{(SDE)}\qquad \mathrm{d}x = -2\dot{\sigma}(t)\sigma(t)\nabla_x\log p_t(x)dt + \sqrt{2\dot{\sigma}(t)\sigma(t)}\mathrm{d}W_t \numberthis\label{eq: diff-sde}\\
       &\text{(ODE)}\qquad \mathrm{d}x = -\dot{\sigma}(t)\sigma(t)\nabla_x\log p_t(x)dt \numberthis\label{eq: diff-ode}
\end{align*}
where $\nabla_x\log p_{t}(x)$ in the drift term is the score of intermediate distribution at time $t$. W.l.o.g we set $\sigma(t)=t$ in the remaining text, as in \cite{Karras2022ElucidatingTD}. Both processes progressively recover $p_0$ from the prior distribution $p_T$ while sharing the same time-dependent distribution $p_t$. In practice, we train a neural network $s_\theta(x, t)$ to estimate the score field $\nabla_x\log p_t(x)$ by minimizing the denoising score-matching loss~\cite{Vincent2011ACB}. We then substitute the score $\nabla_x\log p_{t}(x)$ with $s_\theta(x, t)$ in the drift term of above backward SDE~(\Eqref{eq: diff-sde})/ODE~(\Eqref{eq: diff-ode}) for sampling.

Recent work inspired by electrostatics has not only challenged but also integrated diffusion models, notably PFGM/PFGM++, enhances performance in both image and antibody generation~\cite{Xu2022PoissonFG, Xu2023PFGMUT, antipfgm}. They interpret data as electric charges in an augmented space, and the generative processes involve the simulations of differential equations defined by electric field lines. Similar to diffusion models, PFGMs train a neural network to approximate the electric field in the augmented space.
